% Mathematical appendix supporting the theoretical and simulation chapters.
\chapter{Definitions and Supporting Proofs}
\label{app:mathematical-proofs}

This appendix collects the definitions and mathematical arguments used from
Chapter~\ref{chap:introduction} through the estimation and inference
chapters.  Its purpose is to
make the main text self-contained without interrupting the exposition with
long proofs.  A definition fixes the meaning of a term and therefore is not a
statement that can be proved.  What can be proved are the properties and
consequences that follow from those definitions.  Accordingly,
Section~\ref{sec:app-definitions} records the terminology, and the remaining
sections prove the results invoked in the thesis.  The presentation follows
standard Hilbert-space and inverse-problem arguments; see
\textcite{lin_mit_18102_2021}, \textcite{hsing_eubank_2015}, and
\textcite{engl_hanke_neubauer_1996} for broader treatments.

Throughout, $\mathcal H$ is a real separable Hilbert space with inner product
$\langle\cdot,\cdot\rangle$ and norm $\|\cdot\|$.  In the functional
regression model, $\mathcal H=L^2(\mathcal I)$ for a compact interval
$\mathcal I$.  All equalities between elements of $L^2(\mathcal I)$ are
understood up to sets of Lebesgue measure zero.
When the covariance operator $K$ is not injective, slope parameters are
reported through their minimum-norm representatives in the identified space
\[
    \mathcal H_{\mathrm{id}}
    =\operatorname{Null}(K)^\perp
    =\overline{\operatorname{Range}(K)}.
\]
Statements about a globally identified slope implicitly restrict the
parameter space to $\mathcal H_{\mathrm{id}}$, unless injectivity of $K$ is
assumed explicitly.

\section{Definitions and notation}
\label{sec:app-definitions}

\begin{definition}[Inner product, norm, and orthogonality]
An \emph{inner product} on a real vector space is a symmetric bilinear map
$\langle\cdot,\cdot\rangle$ satisfying $\langle f,f\rangle>0$ for every
$f\neq0$.  Its induced norm is $\|f\|=\langle f,f\rangle^{1/2}$.  Elements
$f$ and $g$ are \emph{orthogonal}, written $f\perp g$, when
$\langle f,g\rangle=0$.
\end{definition}

\begin{definition}[Hilbert space and orthonormal basis]
A \emph{Hilbert space} is an inner-product space that is complete in its
induced norm.  A sequence $\{e_j\}_{j\geq1}$ is \emph{orthonormal} if
$\langle e_j,e_k\rangle$ equals one for $j=k$ and zero otherwise.  It is an
\emph{orthonormal basis} if the closure of its linear span is the whole
space.
\end{definition}

\begin{definition}[$\ell^p$ spaces, conjugate exponents, and dual balls]
For $1\leq p<\infty$, the sequence space $\ell^p$ consists of all real
sequences $x=(x_j)_{j\geq1}$ for which
\[
    \|x\|_p=\left(\sum_{j\geq1}|x_j|^p\right)^{1/p}<\infty.
\]
The space $\ell^\infty$ consists of bounded sequences and has norm
$\|x\|_\infty=\sup_j|x_j|$.  Exponents $p,q\in[1,\infty]$ are
\emph{H\"older conjugates} when $p^{-1}+q^{-1}=1$, using the convention
$\infty^{-1}=0$.  The continuous dual $E^*$ of a normed space $E$ is the
space of its continuous linear functionals, equipped with
\[
    \|\varphi\|_*=\sup_{\|x\|_E\leq1}|\varphi(x)|.
\]
In $\mathbb R^d$, write $B_p^d=\{x:\|x\|_p\leq1\}$.  Its polar under the
standard pairing is
\[
    (B_p^d)^\circ
    =\left\{y:\sup_{x\in B_p^d}|y^{\mathsf T}x|\leq1\right\}.
\]
The norm-induced distance is $d_p(x,z)=\|x-z\|_p$; its radius-$r$ ball
centred at $z$ is $z+rB_p^d$.
\end{definition}

\begin{definition}[Linear operators]
For a linear operator $A:\mathcal H\to\mathcal H$, its null space and range
are
\[
    \operatorname{Null}(A)=\{f:Af=0\},
    \qquad
    \operatorname{Range}(A)=\{Af:f\in\mathcal H\}.
\]
The operator is \emph{injective} when $\operatorname{Null}(A)=\{0\}$ and is
\emph{bounded} when its operator norm is finite, where
\[
    \|A\|_{\mathrm{op}}
    =\sup_{\|f\|\leq1}\|Af\|.
\]
A bounded operator is \emph{self-adjoint} if
$\langle Af,g\rangle=\langle f,Ag\rangle$ for all $f,g$, and
\emph{non-negative} if $\langle Af,f\rangle\geq0$ for every $f$.  An
\emph{eigenpair} $(\lambda,v)$ satisfies $Av=\lambda v$ with $v\neq0$.
\end{definition}

\begin{definition}[Compact and trace-class operators]
A bounded operator $A$ is \emph{compact} if it maps every bounded sequence to
a sequence whose image has a convergent subsequence.  A non-negative
self-adjoint operator is \emph{trace class} if, for one and hence every
orthonormal basis $\{e_j\}$,
\(
  \operatorname{tr}(A)=\sum_j\langle Ae_j,e_j\rangle<\infty.
\)
Every trace-class operator is compact.
\end{definition}

\begin{definition}[Rank-one and covariance operators]
For $u,v\in\mathcal H$, the rank-one operator $u\otimes v$ is defined by
$(u\otimes v)f=\langle v,f\rangle u$.  If $X$ is a random element with mean
$\mu_X$ and $\mathbb E\|X\|^2<\infty$, its covariance operator is
\[
    K=\mathbb E\{(X-\mu_X)\otimes(X-\mu_X)\}.
\]
\end{definition}

\begin{definition}[Well-posedness, ill-posedness, and regularisation]
An inverse problem is \emph{well posed} in the sense relevant here if a
solution exists, is unique on the chosen parameter space, and depends
continuously on the observed moments.  Failure of one of these conditions
makes the problem \emph{ill posed}.  A \emph{regularisation method} replaces
an unstable inverse by a family of stable approximations indexed by a tuning
parameter, truncation level, or stopping index.
\end{definition}

\begin{definition}[Condition number and Krylov subspace]
For a tall matrix $A\in\mathbb R^{p\times m}$ with $p\geq m$, its spectral
condition number is
\[
    \kappa_2(A)=\frac{\sigma_{\max}(A)}{\sigma_{\min}(A)}
\]
when $A$ has full column rank, and $\kappa_2(A)=\infty$ otherwise.  For a
symmetric positive-definite square matrix this reduces to
$\kappa_2(A)=\lambda_{\max}(A)/\lambda_{\min}(A)$.  For an operator $A$, a
right-hand side $r$, and an integer $m\geq1$, the order-$m$ Krylov subspace is
\[
    \mathcal K_m(A,r)
    =\operatorname{span}\{r,Ar,\ldots,A^{m-1}r\}.
\]
\end{definition}

\begin{definition}[Sampling operator, adjoint, and moment residual]
For centred functional observations $X_1^c,\ldots,X_n^c$, the sampling
operator is $(T_nb)_i=\langle X_i^c,b\rangle$.  Relative to the sample inner
product $\langle a,c\rangle_n=n^{-1}a^\top c$, its adjoint is
$T_n^*a=n^{-1}\sum_i a_iX_i^c$.  Hence
$\widehat K=T_n^*T_n$, $\widehat r=T_n^*y$, and the empirical moment residual
at $b$ is $e(b)=\widehat r-\widehat Kb=T_n^*(y-T_nb)$.
\end{definition}

\begin{definition}[PLS1, NIPALS, and APLS]
\emph{PLS1} denotes partial least squares with one scalar response.
\emph{NIPALS} is a classical sequential PLS algorithm that constructs a
response-covariant score, estimates predictor and response loadings, and
deflates both residual objects after each component.  In the present
functional setting, \emph{APLS} denotes the alternative representation that
minimises response least squares over
$\mathcal K_m(\widehat K,\widehat r)$.  Under the full-rank conditions stated
below, these are two representations of the same fixed-component standard
functional PLS fit.
\end{definition}

\begin{definition}[Arnoldi basis and conjugacy]
An Arnoldi basis of $\mathcal K_m(A,r)$ is an orthonormal sequence generated
by repeatedly applying $A$ to the latest basis vector and orthogonalising the
result against the existing basis.  Non-zero vectors $p_0,\ldots,p_k$ are
\emph{$A$-conjugate} if $\langle Ap_j,p_\ell\rangle=0$ whenever
$j\neq\ell$.
\end{definition}

\begin{definition}[Stochastic order and weak convergence]
For random quantities $Z_n$, the notation $Z_n=O_{\mathbb P}(1)$ means that
$\{Z_n\}$ is bounded in probability, while $Z_n=o_{\mathbb P}(1)$ means
$Z_n\to0$ in probability.  The notation $Z_n\rightsquigarrow Z$ denotes
convergence in distribution.  These definitions also apply to norms of
random elements in $\mathcal H$.
\end{definition}

\section{Hilbert-space geometry}
\label{sec:app-hilbert-proofs}

\begin{theorem}[H\"older duality for $\ell^p$]
\label{thm:app-lp-duality}
Let $1\leq p<\infty$ and let $q$ be its H\"older conjugate.  If
$x\in\ell^p$ and $y\in\ell^q$, then
\begin{equation}
    \left|\sum_{j\geq1}x_jy_j\right|
    \leq \|x\|_p\|y\|_q.
    \label{eq:app-holder-sequence}
\end{equation}
The map $y\mapsto\varphi_y$, where
$\varphi_y(x)=\sum_{j\geq1}x_jy_j$, is an isometric isomorphism from
$\ell^q$ onto $(\ell^p)^*$.  In finite dimension this implies
\[
    (B_p^d)^\circ=B_q^d,
    \qquad 1\leq p\leq\infty.
\]
Consequently, $\ell^2$ is self-dual and $(B_2^d)^\circ=B_2^d$.
\end{theorem}

\begin{proof}
First suppose $1<p<\infty$.  Young's inequality
$ab\leq a^p/p+b^q/q$ for $a,b\geq0$ gives, after normalising non-zero
$x$ and $y$ by their respective norms,
\[
  \sum_{j\geq1}
  \frac{|x_j|}{\|x\|_p}\frac{|y_j|}{\|y\|_q}
  \leq
  \frac1p\sum_{j\geq1}\frac{|x_j|^p}{\|x\|_p^p}
  +\frac1q\sum_{j\geq1}\frac{|y_j|^q}{\|y\|_q^q}
  =1.
\]
This proves \eqref{eq:app-holder-sequence}; if either sequence is zero the
claim is immediate.  When $p=1$ and $q=\infty$, the same conclusion follows
directly from
$\sum_j|x_jy_j|\leq(\sup_j|y_j|)\sum_j|x_j|$.

H\"older's inequality shows that every $y\in\ell^q$ defines a continuous
functional with $\|\varphi_y\|_*\leq\|y\|_q$.  For
$1<p<\infty$ and $y\neq0$, take $N$ such that
$\sum_{k=1}^{N}|y_k|^q>0$; this holds for all sufficiently large $N$.
Equality is obtained as a limit of the finite-support vectors
\[
  x_j^{(N)}=
  \begin{cases}
  \displaystyle
  \frac{\operatorname{sgn}(y_j)|y_j|^{q-1}}
       {\left(\sum_{k=1}^{N}|y_k|^q\right)^{(q-1)/q}},
       &j\leq N,\\[1.2ex]
  0,&j>N,
  \end{cases}
\]
for which $\|x^{(N)}\|_p=1$ and
$\varphi_y(x^{(N)})=(\sum_{j=1}^{N}|y_j|^q)^{1/q}$.
For $p=1$, choose signed coordinate vectors whose coordinates approach
$\|y\|_\infty$.  Hence $\|\varphi_y\|_*=\|y\|_q$ in both cases.

Conversely, let $\varphi\in(\ell^p)^*$ and put
$y_j=\varphi(e_j)$ for the standard coordinate vectors.  If $p=1$, then
$|y_j|\leq\|\varphi\|_*$, so $y\in\ell^\infty$.  If $1<p<\infty$, apply
$\varphi$ to the preceding norm-one construction based on
$(y_1,\ldots,y_N)$ whenever its $q$th-power sum is positive.  When that
sum is zero, the following inequality is immediate.  Thus, in either case,
\[
   \left(\sum_{j=1}^{N}|y_j|^q\right)^{1/q}
   \leq\|\varphi\|_*
   \qquad\text{for every }N,
\]
and therefore $y\in\ell^q$.  On every finitely supported $x$,
$\varphi(x)=\sum_jx_jy_j$.  Such sequences are dense in $\ell^p$ for
$p<\infty$, so continuity extends the identity to every $x\in\ell^p$.
Thus $y\mapsto\varphi_y$ is onto as well as isometric.

In $\mathbb R^d$, the dual-norm identity says
$\sup_{x\in B_p^d}|y^{\mathsf T}x|=\|y\|_q$, which is exactly
$(B_p^d)^\circ=B_q^d$.  The endpoint $p=\infty$, $q=1$ follows by taking
$x_j=\operatorname{sgn}(y_j)$; this is a finite-dimensional polar identity,
not an assertion that $(\ell^\infty)^*=\ell^1$.  Finally, $p=q=2$ yields
self-duality.  The Hilbert-space version of the same identification is the
Riesz representation theorem proved below.
\end{proof}

\begin{remark}[The $\ell^\infty$ endpoint]
\label{rem:app-linfty-dual}
The finite-dimensional identity $(B_\infty^d)^\circ=B_1^d$ does not extend
to an identification $(\ell^\infty)^*=\ell^1$.  Every $y\in\ell^1$ does
define a continuous functional on $\ell^\infty$ through coordinate pairing,
but the full dual $(\ell^\infty)^*$ is strictly larger.  Establishing that
strict inclusion requires additional functional-analysis machinery and is not
used elsewhere in the thesis; it is stated here only to prevent the
finite-dimensional drawing from being overinterpreted
\parencite{lin_mit_18102_2021}.
\end{remark}

\begin{theorem}[Cauchy--Schwarz inequality]
\label{thm:app-cauchy-schwarz}
For all $f,g\in\mathcal H$,
\[
    |\langle f,g\rangle|\leq\|f\|\,\|g\|.
\]
\end{theorem}

\begin{proof}
If $g=0$, the assertion is immediate.  Otherwise, non-negativity of the
squared norm gives, for every $t\in\mathbb R$,
\[
    0\leq\|f-tg\|^2
    =\|f\|^2-2t\langle f,g\rangle+t^2\|g\|^2.
\]
Choose $t=\langle f,g\rangle/\|g\|^2$.  Substitution yields
$0\leq\|f\|^2-|\langle f,g\rangle|^2/\|g\|^2$, which is equivalent to the
claimed inequality.
\end{proof}

\begin{theorem}[Completeness and separability of \texorpdfstring{$L^2$}{L2}]
\label{thm:app-l2-hilbert}
For a compact interval $\mathcal I$, the space $L^2(\mathcal I)$ with its
usual inner product is a separable Hilbert space.
\end{theorem}

\begin{proof}
It remains to verify completeness and separability.  Let $\{f_n\}$ be Cauchy
in $L^2(\mathcal I)$.  Choose a subsequence $\{f_{n_k}\}$ such that
\[
    \|f_{n_{k+1}}-f_{n_k}\|_2\leq2^{-k}.
\]
Put $g_k=f_{n_{k+1}}-f_{n_k}$.  Since $\mathcal I$ has finite measure,
Cauchy--Schwarz gives
$\|g_k\|_1\leq|\mathcal I|^{1/2}\|g_k\|_2$.  Hence
$\sum_k\int_{\mathcal I}|g_k|<\infty$, so $\sum_k|g_k(t)|<\infty$ for almost
every $t$.  Define $f=f_{n_1}+\sum_{k\geq1}g_k$ at those points.  Minkowski's
inequality and Fatou's lemma imply
\[
 \left\|f-f_{n_m}\right\|_2
 \leq\sum_{k=m}^{\infty}\|g_k\|_2
 \leq\sum_{k=m}^{\infty}2^{-k}\longrightarrow0.
\]
Thus the subsequence converges in $L^2$; because the original sequence is
Cauchy, the whole sequence converges to the same $f$.  This proves
completeness.

For separability, consider step functions taking rational values on finite
partitions whose endpoints are rational points of $\mathcal I$.  This family
is countable.  Simple functions are dense in $L^2$, and measurable sets of
finite measure can be approximated in measure by finite unions of intervals;
their endpoints and the corresponding function values can then be
approximated by rationals.  The stated countable family is therefore dense in
$L^2(\mathcal I)$.
\end{proof}

\begin{theorem}[Orthogonal projection]
\label{thm:app-projection}
Let $M$ be a closed linear subspace of $\mathcal H$.  For every
$f\in\mathcal H$, there exists a unique $P_Mf\in M$ such that
$f-P_Mf\perp M$.  Moreover,
\[
    \|f-P_Mf\|=\inf_{v\in M}\|f-v\|.
\]
If $M=\operatorname{span}\{e_1,\ldots,e_m\}$ for orthonormal vectors, then
$P_Mf=\sum_{j=1}^m\langle f,e_j\rangle e_j$.
\end{theorem}

\begin{proof}
Put $d=\inf_{v\in M}\|f-v\|$ and choose $v_n\in M$ with
$\|f-v_n\|\to d$.  The parallelogram identity and
$(v_n+v_k)/2\in M$ imply
\[
 \|v_n-v_k\|^2
 \leq 2\|f-v_n\|^2+2\|f-v_k\|^2-4d^2,
\]
so $\{v_n\}$ is Cauchy.  Completeness of $\mathcal H$ and closedness of $M$
give a limit $v\in M$ with $\|f-v\|=d$.  For $h\in M$, the function
$t\mapsto\|f-v-th\|^2$ is minimised at zero; its linear term therefore
vanishes, giving $\langle f-v,h\rangle=0$.  If $v'$ has the same property,
then $v-v'\in M$ and $v-v'\perp M$, so $v=v'$.  This proves existence,
orthogonality, and uniqueness.  Pythagoras gives
$\|f-w\|^2=\|f-v\|^2+\|v-w\|^2$ for every $w\in M$, which establishes best
approximation.  In the finite orthonormal case, direct substitution shows
that $f-\sum_{j=1}^m\langle f,e_j\rangle e_j$ is orthogonal to every basis
vector and hence to $M$.
\end{proof}

\begin{corollary}[Basis expansion and Parseval's identity]
\label{cor:app-parseval}
If $\{e_j\}_{j\geq1}$ is an orthonormal basis of $\mathcal H$, then, for
every $f\in\mathcal H$,
\[
 f=\sum_{j=1}^{\infty}\langle f,e_j\rangle e_j,
 \qquad
 \|f\|^2=\sum_{j=1}^{\infty}|\langle f,e_j\rangle|^2,
\]
where the series converges in norm.
\end{corollary}

\begin{proof}
Let $P_mf=\sum_{j=1}^m\langle f,e_j\rangle e_j$.  Because the linear span
of the basis is dense, the best-approximation property in
Theorem~\ref{thm:app-projection} implies $\|f-P_mf\|\to0$.  Orthogonality and
Pythagoras give
\[
    \|f\|^2=\|f-P_mf\|^2
       +\sum_{j=1}^m|\langle f,e_j\rangle|^2.
\]
Letting $m\to\infty$ proves both assertions.
\end{proof}

\begin{theorem}[Riesz representation]
\label{thm:app-riesz}
Every continuous linear functional $\ell:\mathcal H\to\mathbb R$ has a
unique representative $b\in\mathcal H$ such that
$\ell(f)=\langle f,b\rangle$ for all $f\in\mathcal H$.
\end{theorem}

\begin{proof}
If $\ell=0$, take $b=0$.  Otherwise, $M=\operatorname{Null}(\ell)$ is a
closed subspace.  Select $z\notin M$ and put $u=z-P_Mz$.  Then $u\neq0$ and
$u\perp M$.  For arbitrary $f$, the element
\[
    f-\frac{\ell(f)}{\ell(u)}u
\]
belongs to $M$.  Consequently,
$f=m+\{\ell(f)/\ell(u)\}u$ for some $m\in M$.  With
$b=\ell(u)u/\|u\|^2$, orthogonality gives
$\langle f,b\rangle=\ell(f)$.  If both $b$ and $b'$ represent $\ell$, then
$\langle f,b-b'\rangle=0$ for every $f$; taking $f=b-b'$ proves $b=b'$.
\end{proof}

\section{Random functions and covariance operators}
\label{sec:app-covariance-proofs}

\begin{proposition}[Existence of the Hilbert-space mean]
\label{prop:app-functional-mean}
If $X$ is a random element of $\mathcal H$ with
$\mathbb E\|X\|^2<\infty$, there is a unique $\mu_X\in\mathcal H$ satisfying
\[
    \langle\mu_X,f\rangle=\mathbb E\langle X,f\rangle
    \quad\text{for every }f\in\mathcal H.
\]
Moreover, $\|\mu_X\|\leq(\mathbb E\|X\|^2)^{1/2}$.
\end{proposition}

\begin{proof}
The map $\ell(f)=\mathbb E\langle X,f\rangle$ is linear.  By
Theorem~\ref{thm:app-cauchy-schwarz} and Cauchy--Schwarz for scalar random
variables,
\[
 |\ell(f)|\leq\mathbb E(\|X\|\|f\|)
 \leq(\mathbb E\|X\|^2)^{1/2}\|f\|.
\]
It is therefore continuous.  Theorem~\ref{thm:app-riesz} gives the unique
representative $\mu_X$, and the displayed bound gives the norm inequality.
\end{proof}

\begin{theorem}[Properties of a covariance operator]
\label{thm:app-covariance-properties}
Let $Z=X-\mu_X$ and suppose $\mathbb E\|Z\|^2<\infty$.  The operator
\[
    Kf=\mathbb E\{\langle Z,f\rangle Z\}
\]
is bounded, self-adjoint, non-negative, trace class, and compact.  In
particular,
\[
    \|K\|_{\mathrm{op}}\leq\mathbb E\|Z\|^2,
    \qquad
    \operatorname{tr}(K)=\mathbb E\|Z\|^2.
\]
\end{theorem}

\begin{proof}
For $f,g\in\mathcal H$,
\[
 \langle Kf,g\rangle
 =\mathbb E\{\langle Z,f\rangle\langle Z,g\rangle\}.
\]
The expression is symmetric in $f$ and $g$, proving self-adjointness.  With
$g=f$ it is non-negative.  Moreover,
\[
 |\langle Kf,g\rangle|
 \leq\mathbb E(\|Z\|^2)\|f\|\|g\|,
\]
which proves boundedness and the operator-norm bound.

For any orthonormal basis $\{e_j\}$, Tonelli's theorem and
Corollary~\ref{cor:app-parseval} yield
\[
 \sum_{j=1}^{\infty}\langle Ke_j,e_j\rangle
 =\mathbb E\sum_{j=1}^{\infty}|\langle Z,e_j\rangle|^2
 =\mathbb E\|Z\|^2<\infty.
\]
Thus, $K$ is trace class and has the asserted trace.

For completeness, compactness can also be seen directly.  Let $P_m$ project
onto
\[
    M_m=\operatorname{span}\{e_1,\ldots,e_m\},
\]
and set $K_m=P_mKP_m$, which has finite rank.  For unit vectors $f,g$,
Cauchy--Schwarz gives
\[
 |\langle (I-P_m)Kf,g\rangle|
 \leq
 \{\mathbb E\|Z\|^2\}^{1/2}
 \{\mathbb E\|(I-P_m)Z\|^2\}^{1/2}.
\]
The final expectation tends to zero by Parseval and dominated convergence.
The same bound applies to $P_mK(I-P_m)$, so
$\|K-K_m\|_{\mathrm{op}}\to0$.  An operator-norm limit of finite-rank
operators is compact.
\end{proof}

\begin{proposition}[Covariance-kernel representation]
\label{prop:app-covariance-kernel}
Suppose that $Z=X-\mu_X$ has pointwise second moments and that the
integrability conditions for Fubini's theorem hold.  If
$c(s,t)=\mathbb E\{Z(s)Z(t)\}$, then, for almost every $s$,
\[
    (Kf)(s)=\int_{\mathcal I}c(s,t)f(t)\,\mathrm dt.
\]
\end{proposition}

\begin{proof}
Using the definition of $K$ and then Fubini's theorem,
\begin{align*}
 (Kf)(s)
 &=\mathbb E\!\left[Z(s)
      \int_{\mathcal I}Z(t)f(t)\,\mathrm dt\right] \\
 &=\int_{\mathcal I}\mathbb E\{Z(s)Z(t)\}f(t)\,\mathrm dt
 =\int_{\mathcal I}c(s,t)f(t)\,\mathrm dt.
\end{align*}
The stated integrability assumptions justify interchanging expectation and
integration.
\end{proof}

\begin{theorem}[Spectral representation]
\label{thm:app-spectral-representation}
Let $K$ be compact, self-adjoint, and non-negative.  Its positive eigenvalues
can be indexed as $\{\lambda_j:j\in\mathcal J_K\}$ in non-increasing order,
where $\mathcal J_K=\{1,\ldots,q\}$ if $K$ has finite rank $q$ and
$\mathcal J_K=\mathbb N$ if it has infinite rank (with
$\mathcal J_K=\varnothing$ when $K=0$).  There are corresponding
orthonormal eigenfunctions $\{v_j:j\in\mathcal J_K\}$ spanning
$\overline{\operatorname{Range}(K)}$ such that
\[
    Kf=\sum_{j\in\mathcal J_K}
        \lambda_j\langle f,v_j\rangle v_j.
\]
If $K$ has infinite rank, then $\lambda_j\to0$.
\end{theorem}

\begin{proof}
If $K=0$, the result is immediate.  Otherwise, define the maximal Rayleigh
quotient
$\lambda_1=\sup_{\|f\|=1}\langle Kf,f\rangle>0$.  Take a maximising
sequence.  The unit ball is weakly sequentially compact, and compactness of
$K$ converts a weakly convergent subsequence $f_n\rightharpoonup v_1$ into
strong convergence $Kf_n\to Kv_1$.  Hence
$\langle Kf_n,f_n\rangle\to\langle Kv_1,v_1\rangle=\lambda_1$.  Since
$\langle Kv_1,v_1\rangle\leq\lambda_1\|v_1\|^2$ and $\|v_1\|\leq1$, it
follows that $\|v_1\|=1$.  Considering the Rayleigh quotient along
$v_1+th$ for $h\perp v_1$ shows that $\langle Kv_1,h\rangle=0$; therefore
$Kv_1=\lambda_1v_1$.

Self-adjointness makes $v_1^\perp$ invariant under $K$.  Repeating the same
argument on successive orthogonal complements produces orthonormal
eigenvectors with non-increasing eigenvalues until the restriction of $K$
is zero.  In the infinite-rank case, if infinitely many eigenvalues were
bounded below by some $c>0$, then the sequence
$Kv_j=\lambda_jv_j$ would have pairwise distances at least $\sqrt{2}c$ and
hence no convergent subsequence, contradicting compactness.  Thus
$\lambda_j\to0$ in that case.

Finally, if a vector orthogonal to every constructed $v_j$ had non-zero
image under $K$, the Rayleigh-quotient argument on that invariant orthogonal
complement would produce a unit eigenvector with a positive eigenvalue
$\rho$.  In a finite construction this contradicts termination.  In an
infinite construction that vector belongs to every preceding orthogonal
complement, so maximality of each Rayleigh quotient gives
$\lambda_j\geq\rho$ for every $j$, contradicting $\lambda_j\to0$.
Hence $K$ is zero on the remaining complement.  This gives the expansion, and the eigenvectors
with positive eigenvalues span $\overline{\operatorname{Range}(K)}$.
\end{proof}

\begin{proposition}[Karhunen--Lo\`eve expansion and optimality]
\label{prop:app-karhunen-loeve}
Under the assumptions of Theorem~\ref{thm:app-covariance-properties}, let
$(\lambda_j,v_j)$ be the positive eigenpairs of $K$ and set
$\xi_j=\langle X-\mu_X,v_j\rangle$.  Then
\[
 X=\mu_X+\sum_{j\in\mathcal J_K}\xi_jv_j
 \quad\text{in mean square},
\]
$\mathbb E\xi_j=0$, and
$\mathbb E(\xi_j\xi_k)=\lambda_j$ for $j=k$ and zero otherwise.  For
each $m\leq\operatorname{rank}(K)$, among all $m$-dimensional subspaces,
$\operatorname{span}\{v_1,\ldots,v_m\}$ minimises the expected squared
reconstruction error.  If $K$ has finite rank, the full expansion is the
case $m=\operatorname{rank}(K)$.
\end{proposition}

\begin{proof}
Centredness gives $\mathbb E\xi_j=0$, while
\[
 \mathbb E(\xi_j\xi_k)
 =\langle Kv_j,v_k\rangle
 =\lambda_j\langle v_j,v_k\rangle.
\]
Every component of $Z=X-\mu_X$ in $\operatorname{Null}(K)$ has zero second
moment because, for $h\in\operatorname{Null}(K)$,
\[
    \mathbb E\langle Z,h\rangle^2=\langle Kh,h\rangle=0.
\]
For each $m\leq\operatorname{rank}(K)$, the trace identity therefore implies
\[
 \mathbb E\left\|Z-\sum_{j=1}^m\xi_jv_j\right\|^2
 =\operatorname{tr}(K)-\sum_{j=1}^m\lambda_j.
\]
If $K$ has infinite rank, the right-hand side tends to zero as
$m\to\infty$; if it has finite rank $q$, it equals zero at $m=q$.  This
proves the asserted mean-square expansion.

For any orthonormal $u_1,\ldots,u_m$, the expected reconstruction error is
\[
 \operatorname{tr}(K)-\sum_{\ell=1}^m\langle Ku_\ell,u_\ell\rangle.
\]
Writing each $u_\ell$ in the eigenbasis gives
$\sum_\ell\langle Ku_\ell,u_\ell\rangle=\sum_j\lambda_ja_j$, where
$0\leq a_j\leq1$ and $\sum_j a_j\leq m$.  Since the eigenvalues are ordered,
this sum cannot exceed $\sum_{j=1}^m\lambda_j$.  Equality is attained by
$u_\ell=v_\ell$, proving optimality.
\end{proof}

\begin{proposition}[Rank of the empirical covariance]
\label{prop:app-empirical-rank}
For observations $X_1,\ldots,X_n$ and a fixed population mean $\mu_X$, the
population-centred empirical covariance operator
\[
 \widehat K_{\mu}f=\frac1n\sum_{i=1}^n
 \langle X_i-\mu_X,f\rangle(X_i-\mu_X)
\]
has rank at most $n$.  The sample-centred empirical covariance operator
\[
 \widehat K_{\overline X}f=\frac1n\sum_{i=1}^n
 \langle X_i-\overline X,f\rangle(X_i-\overline X)
\]
has rank at most $n-1$.
\end{proposition}

\begin{proof}
The range of $\widehat K_{\mu}$ is contained in the span of the $n$ curves
$X_i-\mu_X$, so its rank is at most $n$.  The range of
$\widehat K_{\overline X}$ is contained in the span of the $n$ curves
$X_i-\overline X$.  These curves sum to zero, so their span has dimension at
most $n-1$.  The rank of the corresponding operator cannot exceed that
dimension.
\end{proof}

\section{The functional normal equation and inverse problem}
\label{sec:app-regression-proofs}

\begin{proposition}[Centring and the normal equation]
\label{prop:app-normal-equation}
Suppose
\[
 Y=\alpha+\langle X,\beta\rangle+\varepsilon,
 \qquad \mathbb E(\varepsilon\mid X)=0,
\]
with finite second moments.  Then
$\alpha=\mu_Y-\langle\mu_X,\beta\rangle$, and the centred variables satisfy
$Y^c=\langle X^c,\beta\rangle+\varepsilon$.  If
$r=\mathbb E(Y^cX^c)$ and $K$ is the covariance operator of $X$, then
$K\beta=r$.
\end{proposition}

\begin{proof}
Taking expectations and using
$\mathbb E\varepsilon=\mathbb E\{\mathbb E(\varepsilon\mid X)\}=0$ gives
the formula for $\alpha$ and therefore the centred model.  For arbitrary
$f\in\mathcal H$,
\begin{align*}
 \langle r,f\rangle
 &=\mathbb E\{Y^c\langle X^c,f\rangle\} \\
 &=\mathbb E\{\langle X^c,\beta\rangle\langle X^c,f\rangle\}
   +\mathbb E\{\varepsilon\langle X^c,f\rangle\}.
\end{align*}
The last expectation is zero by iterated expectations.  The first equals
$\langle K\beta,f\rangle$.  Since the equality holds for every $f$,
$r=K\beta$.
\end{proof}

\begin{proposition}[Excess prediction risk]
\label{prop:app-risk-identity}
For $R(b)=\mathbb E\{Y^c-\langle X^c,b\rangle\}^2$ under the correctly
specified model,
\[
 R(b)-R(\beta)
 =\langle K(b-\beta),b-\beta\rangle
 =\|K^{1/2}(b-\beta)\|^2.
\]
\end{proposition}

\begin{proof}
Let $h=b-\beta$.  The prediction residual at $b$ is
$\varepsilon-\langle X^c,h\rangle$.  Expanding the square gives
\[
 R(b)=\mathbb E\varepsilon^2
 -2\mathbb E\{\varepsilon\langle X^c,h\rangle\}
 +\mathbb E\langle X^c,h\rangle^2.
\]
The middle term is zero by the conditional mean restriction, and the last
term is $\langle Kh,h\rangle$.  Since $R(\beta)=\mathbb E\varepsilon^2$, the
first equality follows.  The spectral representation defines
$K^{1/2}f=\sum_j\lambda_j^{1/2}\langle f,v_j\rangle v_j$ and gives
$\langle Kh,h\rangle=\|K^{1/2}h\|^2$.
\end{proof}

\begin{proposition}[Variance identities for the discrete simulation]
\label{prop:app-simulation-variance}
Let
\[
 X=\sum_{j=1}^{J}\sqrt{\lambda_j}Z_j\phi_j,
 \qquad
 Y=\langle X,\beta\rangle_T+\varepsilon,
\]
where the $Z_j$ are independent with mean zero and variance one,
$\varepsilon$ is independent of them with mean zero and variance $\sigma^2$,
and $\langle\cdot,\cdot\rangle_T$ is the discrete inner product.  Then
\begin{align*}
 \mathbb E\|X\|_T^2
 &=\sum_{j=1}^{J}\lambda_j\|\phi_j\|_T^2,\\
 \operatorname{Var}\{\langle X,\beta\rangle_T\}
 &=\sum_{j=1}^{J}\lambda_j
   \langle\phi_j,\beta\rangle_T^2,\\
 \operatorname{Var}(Y)
 &=\sum_{j=1}^{J}\lambda_j
   \langle\phi_j,\beta\rangle_T^2+\sigma^2.
\end{align*}
\end{proposition}

\begin{proof}
Write $c_j=\langle\phi_j,\beta\rangle_T$.  Bilinearity gives
\[
 \langle X,\beta\rangle_T
 =\sum_{j=1}^{J}\sqrt{\lambda_j}Z_jc_j.
\]
All cross-products have expectation zero because the scores are independent
and centred, while $\mathbb E Z_j^2=1$.  Therefore the variance of this sum is
$\sum_j\lambda_jc_j^2$.  Similarly,
\[
 \mathbb E\|X\|_T^2
 =\sum_{j,k}\sqrt{\lambda_j\lambda_k}\,
   \mathbb E(Z_jZ_k)\langle\phi_j,\phi_k\rangle_T
 =\sum_j\lambda_j\|\phi_j\|_T^2.
\]
Finally, independence of $\varepsilon$ and the scores makes the covariance
between the noiseless response and the error zero, so their variances add.
No discrete orthogonality assumption on the $\phi_j$ is required.
\end{proof}

\begin{proposition}[Empirical bias--variance identity]
\label{prop:app-mc-decomposition}
For arbitrary $b_1,\ldots,b_R,\beta$ in a real Hilbert space, let
$\overline b=R^{-1}\sum_{r=1}^{R}b_r$.  Then
\[
 \frac1R\sum_{r=1}^{R}\|b_r-\beta\|^2
 =\|\overline b-\beta\|^2
  +\frac1R\sum_{r=1}^{R}\|b_r-\overline b\|^2.
\]
\end{proposition}

\begin{proof}
For every $r$,
\[
 b_r-\beta=(b_r-\overline b)+(\overline b-\beta).
\]
Expanding the squared norm and averaging yields the two displayed squared
terms plus
\[
 \frac{2}{R}\sum_{r=1}^{R}
 \langle b_r-\overline b,\overline b-\beta\rangle.
\]
The cross-term is zero because
$\sum_r(b_r-\overline b)=0$.  The argument applies in particular to the
grid inner product used in the simulation and explains why the empirical
variance uses divisor $R$ when an exact decomposition of the Monte Carlo mean
squared error is desired.
\end{proof}

\begin{proposition}[Identification and the minimum-norm target]
\label{prop:app-identification}
Two slopes $\beta$ and $\widetilde\beta$ generate the same centred
conditional mean if and only if
$\widetilde\beta-\beta\in\operatorname{Null}(K)$.  In the uncentred model,
if $h=\widetilde\beta-\beta\in\operatorname{Null}(K)$, the parameter pairs
$(\alpha,\beta)$ and
$(\alpha-\langle\mu_X,h\rangle,\widetilde\beta)$ generate the same
conditional mean.  Hence the slope is fully identified exactly when $K$ is
injective.  Otherwise, every identified equivalence class has a unique
minimum-norm representative in
\[
 \operatorname{Null}(K)^\perp
 =\overline{\operatorname{Range}(K)}.
\]
\end{proposition}

\begin{proof}
For $h=\widetilde\beta-\beta$,
\[
 \mathbb E\langle X^c,h\rangle^2=\langle Kh,h\rangle.
\]
If $Kh=0$, the non-negative random variable on the left has expectation
zero, so $\langle X^c,h\rangle=0$ almost surely.  Conversely, this almost
sure equality implies
\[
    Kh=\mathbb E\{\langle X^c,h\rangle X^c\}=0.
\]
This proves the centred equivalence and the injectivity statement.  If
$h\in\operatorname{Null}(K)$, then
$\langle X,h\rangle=\langle\mu_X,h\rangle$ almost surely.  Consequently,
\[
 \alpha-\langle\mu_X,h\rangle
 +\langle X,\beta+h\rangle
 =\alpha+\langle X,\beta\rangle
 \quad\text{almost surely},
\]
which proves the uncentred assertion.

Every $b\in\mathcal H$ has the orthogonal decomposition
$b=b_0+h$ with $b_0\in\operatorname{Null}(K)^\perp$ and
$h\in\operatorname{Null}(K)$.  All elements $b_0+h$ give the same
conditional mean, while
$\|b_0+h\|^2=\|b_0\|^2+\|h\|^2$ is uniquely minimised by $h=0$.  Finally,
for every bounded operator,
$\operatorname{Range}(K)^\perp=\operatorname{Null}(K^*)$; self-adjointness
and taking orthogonal complements give
$\overline{\operatorname{Range}(K)}=\operatorname{Null}(K)^\perp$.
\end{proof}

\begin{theorem}[Picard condition and instability of inversion]
\label{thm:app-picard}
Let $K$ be compact, self-adjoint, and non-negative, with positive eigenpairs
$(\lambda_j,v_j)$.  The equation $K\beta=r$ has a solution if and only if
$r\perp\operatorname{Null}(K)$ and
\[
    \sum_{j:\lambda_j>0}
       \frac{|\langle r,v_j\rangle|^2}{\lambda_j^2}<\infty.
\]
When these conditions hold, the unique minimum-norm solution is
\[
 \beta^\dagger
 =\sum_{j:\lambda_j>0}
   \frac{\langle r,v_j\rangle}{\lambda_j}v_j.
\]
If $K$ has infinite rank, the restriction of the minimum-norm inverse to
$\operatorname{Range}(K)$,
$K^\dagger|_{\operatorname{Range}(K)}:\operatorname{Range}(K)
\to\operatorname{Null}(K)^\perp$, is unbounded in the $\mathcal H$ norm.
\end{theorem}

\begin{proof}
If $K\beta=r$, then $r$ lies in
$\overline{\operatorname{Range}(K)}=\operatorname{Null}(K)^\perp$ and
$\langle r,v_j\rangle=\lambda_j\langle\beta,v_j\rangle$.  Parseval therefore
implies the summability condition.  Conversely, the displayed condition
makes the proposed series for $\beta^\dagger$ square summable.  Applying $K$
term by term yields the component of $r$ in
$\operatorname{Null}(K)^\perp$, which equals $r$ by the first condition.
All other solutions differ from $\beta^\dagger$ by an element of
$\operatorname{Null}(K)$, so Proposition~\ref{prop:app-identification} gives
the minimum-norm claim.

For the final statement, $v_j\in\operatorname{Range}(K)$ because
$v_j=K(v_j/\lambda_j)$.  Yet
\[
 \frac{\|K^\dagger v_j\|}{\|v_j\|}
 =\frac1{\lambda_j}\longrightarrow\infty
\]
as $\lambda_j\to0$.  Thus no finite operator-norm bound for $K^\dagger$
exists on $\operatorname{Range}(K)$.
\end{proof}

\section{Regularised and Krylov solutions}
\label{sec:app-regularisation-proofs}

\begin{proposition}[Tikhonov solution]
\label{prop:app-tikhonov}
For $\rho>0$, the criterion $R(b)+\rho\|b\|^2$ has the unique minimiser
\[
    \beta_\rho=(K+\rho I)^{-1}r.
\]
If $r\perp\operatorname{Null}(K)$, its spectral form is
\[
    \beta_\rho
    =\sum_{j:\lambda_j>0}
      \frac{\langle r,v_j\rangle}{\lambda_j+\rho}v_j.
\]
\end{proposition}

\begin{proof}
Apart from a constant independent of $b$, the population risk equals
$\langle Kb,b\rangle-2\langle r,b\rangle$.  The spectral representation
shows directly that $K+\rho I$ is invertible and that
$\|(K+\rho I)^{-1}\|_{\mathrm{op}}\leq1/\rho$.  The directional derivative
of the penalised criterion at $b$ in direction $h$ is
\[
 2\langle (K+\rho I)b-r,h\rangle.
\]
It vanishes for every $h$ exactly when $(K+\rho I)b=r$.  Since
$\langle (K+\rho I)h,h\rangle\geq\rho\|h\|^2$, the criterion is strictly
convex and the solution is unique.  Expansion in the eigenbasis produces the
stated filter.
\end{proof}

\begin{proposition}[FPCR restriction]
\label{prop:app-fpcr}
Let $V_m=\operatorname{span}\{v_1,\ldots,v_m\}$ with
$\lambda_1,\ldots,\lambda_m>0$.  The minimiser of $R(b)$ over $V_m$ is
\[
    \beta_m^{\mathrm{FPCR}}
    =\sum_{j=1}^m
      \frac{\langle r,v_j\rangle}{\lambda_j}v_j.
\]
The empirical formula follows after replacing $(K,r)$ by
$(\widehat K,\widehat r)$.
\end{proposition}

\begin{proof}
Write $b=\sum_{j=1}^m b_jv_j$.  Up to an additive constant,
\[
 R(b)=\sum_{j=1}^m
 \{\lambda_jb_j^2-2b_j\langle r,v_j\rangle\}.
\]
Each term is a strictly convex quadratic whose minimiser is
$b_j=\langle r,v_j\rangle/\lambda_j$.  Combining the coefficients proves
the result.
\end{proof}

\begin{proposition}[Polynomial form of Krylov estimators]
\label{prop:app-krylov-polynomial}
Every $\beta_m\in\mathcal K_m(K,r)$ can be written as
$\beta_m=p_{m-1}(K)r$ for a polynomial of degree at most $m-1$.  Its moment
residual satisfies
\[
 r-K\beta_m=\pi_m(K)r,
 \qquad \pi_m(\lambda)=1-\lambda p_{m-1}(\lambda).
\]
\end{proposition}

\begin{proof}
Membership in $\mathcal K_m(K,r)$ means that there are coefficients
$a_0,\ldots,a_{m-1}$ with
$\beta_m=\sum_{j=0}^{m-1}a_jK^jr$.  Taking
$p_{m-1}(\lambda)=\sum_{j=0}^{m-1}a_j\lambda^j$ gives the first assertion.
Then
$r-K\beta_m=\{I-Kp_{m-1}(K)\}r=\pi_m(K)r$ by direct algebra.
\end{proof}

\begin{proposition}[Raw and Arnoldi Krylov spaces]
\label{prop:app-arnoldi-span}
In exact arithmetic, if Arnoldi orthogonalisation does not break down before
step $m$, its vectors $q_1,\ldots,q_m$ form an orthonormal basis of
$\mathcal K_m(K,r)$.  Consequently, a uniquely defined restricted
least-squares solution is invariant to whether the raw power basis or the
Arnoldi basis is used.
\end{proposition}

\begin{proof}
The first Arnoldi vector is $q_1=r/\|r\|$ and hence spans
$\mathcal K_1(K,r)$.  Suppose inductively that
$\operatorname{span}\{q_1,\ldots,q_j\}=\mathcal K_j(K,r)$.  Arnoldi forms
$Kq_j$ and removes its projections onto the existing vectors.  The remainder
therefore belongs to $\mathcal K_{j+1}(K,r)$ and is orthogonal to
$\mathcal K_j(K,r)$.  Absence of breakdown means that this remainder is
non-zero, so its normalisation adds one new direction and spans
$\mathcal K_{j+1}(K,r)$.  Induction proves the first claim.  The optimisation
problem depends on the subspace, not on its coordinates; uniqueness therefore
makes its solution basis-invariant.  Finite-precision discrepancies arise
because computed raw powers may lose linear independence, not because the
exact subspaces differ.
\end{proof}

\begin{proposition}[Response-space Krylov characterisation]
\label{prop:app-response-krylov-variational}
Let $T:\mathcal H\to\mathbb R^n$ be bounded, let $K=T^*T$ and $r=T^*y$,
and put $M_m=\mathcal K_m(K,r)$.  If $T$ is injective on $M_m$, the
restricted response least-squares problem
\[
    \min_{b\in M_m}\|y-Tb\|_n^2
\]
has a unique solution $\beta_m^Y$.  Its moment residual
$e_m^Y=r-K\beta_m^Y$ satisfies
\[
    e_m^Y\perp M_m.
\]
Equivalently, $\beta_m^Y$ uniquely minimises
$\Phi(b)=\frac12\langle Kb,b\rangle-\langle r,b\rangle$ over $M_m$.
\end{proposition}

\begin{proof}
Expanding the response criterion and using $K=T^*T$ and $r=T^*y$ gives
\[
 \frac12\|y-Tb\|_n^2
 =\frac12\|y\|_n^2
  +\frac12\langle Kb,b\rangle-\langle r,b\rangle.
\]
The first term is constant in $b$.  For $h\in M_m$, the directional
derivative of the remaining expression is
\[
 \left.\frac{\mathrm d}{\mathrm dt}\Phi(\beta_m^Y+th)\right|_{t=0}
 =\langle K\beta_m^Y-r,h\rangle.
\]
It vanishes for every $h\in M_m$ precisely when $e_m^Y\perp M_m$.
Injectivity of $T$ on the finite-dimensional space gives
$\langle Kh,h\rangle=\|Th\|_n^2>0$ for every non-zero $h\in M_m$.
The criterion is therefore strictly convex on $M_m$, so this stationary point
exists uniquely and is the claimed minimiser.
\end{proof}

\begin{proposition}[NIPALS--APLS equivalence]
\label{prop:app-nipals-apls-equivalence}
Let $T:\mathcal H\to\mathbb R^n$ be bounded, set $K=T^*T$ and $r=T^*y$,
and consider scalar-response PLS with common centring and inner products.
Suppose that the first $m$ score directions generated by the NIPALS-type
score-and-deflation construction are non-zero and that $T$ is injective on
$M_m=\mathcal K_m(K,r)$.  Then the NIPALS score space is $T M_m$, its fitted
response is the orthogonal projection of $y$ onto $T M_m$, and its
reconstructed slope equals the unique APLS solution in
Proposition~\ref{prop:app-response-krylov-variational}.
\end{proposition}

\begin{proof}[Proof sketch]
At the first step, the NIPALS weight is proportional to
$T^*y=r$, so its score lies in $TM_1$.  Predictor deflation at each subsequent
step replaces the current sampling operator $T_{h-1}$ by
$(I-P_{t_h})T_{h-1}$, where $P_{t_h}$ is the response-space orthogonal
projection onto the current score.  Response deflation applies the same
projection to the current response.  Expanding these projections and
repeatedly using $K=T^*T$ gives the standard PLS1 polynomial induction: the
unnormalised $h$th weight has the form
\[
    c_h=\pi_{h-1}(K)r,
\]
where $\pi_{h-1}$ has degree $h-1$.  Its leading coefficient is non-zero
whenever the new score is non-zero.  Consequently,
\[
 \operatorname{span}\{c_1,\ldots,c_h\}
 =\operatorname{span}\{r,Kr,\ldots,K^{h-1}r\}=M_h,
\]
because the change-of-basis matrix between the two displayed sequences is
triangular with non-zero diagonal.  Applying $T$ gives equality of the score
space and $TM_h$.

The predictor deflation makes successive scores mutually orthogonal.  At
each step, the NIPALS response loading is the least-squares coefficient of the
current residual response on the new score.  Hence the sum of the first $m$
fitted score contributions is the orthogonal projection of $y$ onto their
span, namely $TM_m$.  Minimising $\|y-Tb\|_n^2$ over $b\in M_m$ produces the
same projection.  Injectivity on $M_m$ makes the restricted slope unique, so
the reconstructed NIPALS slope and the APLS slope coincide.  This operator
argument is the standard fixed-component equivalence established for
functional PLS by \textcite{delaigle_hall_2012}.
\end{proof}

\begin{proposition}[Response and moment Krylov objectives]
\label{prop:app-response-moment-krylov}
Let $K$ be self-adjoint and non-negative, $r\in\mathcal H$, and
$M_m=\mathcal K_m(K,r)$.  Assume that $\dim(M_m)=m$ and that $K$ is
injective on $M_m$.  Then both restricted quadratic criteria below are
strictly convex.  Define
\[
 \beta_m^Y
 =\arg\min_{b\in M_m}
   \left\{\frac12\langle Kb,b\rangle-\langle r,b\rangle\right\},
 \qquad
 \beta_m^M
 =\arg\min_{b\in M_m}\frac12\|r-Kb\|^2.
\]
These response-space and moment-residual solutions are characterised by
\begin{align*}
    r-K\beta_m^Y&\perp M_m,\\
    r-K\beta_m^M&\perp KM_m.
\end{align*}
With $g_j=K^{j-1}r$ and
$\beta_m=\sum_{j=1}^{m}a_jg_j$, their coefficient systems are respectively
\begin{align*}
 \sum_{k=1}^{m}\langle r,K^{j+k-1}r\rangle a_k
 &=\langle r,K^{j-1}r\rangle,\\
 \sum_{k=1}^{m}\langle r,K^{j+k}r\rangle a_k
 &=\langle r,K^jr\rangle,
 \qquad j=1,\ldots,m.
\end{align*}
Thus, the two methods use the same Krylov space but do not generally have the
same finite-$m$ solution.
\end{proposition}

\begin{proof}
For the response criterion
\[
    \Phi(b)=\frac12\langle Kb,b\rangle-\langle r,b\rangle,
\]
the directional derivative at $b$ in direction $h\in M_m$ is
\[
    \left.\frac{\mathrm d}{\mathrm dt}\Phi(b+th)\right|_{t=0}
    =-\langle r-Kb,h\rangle.
\]
It vanishes for every $h\in M_m$ if and only if
$r-Kb\perp M_m$.  For the moment criterion
\[
    \Psi(b)=\frac12\|r-Kb\|^2,
\]
the directional derivative at $b$ in direction $h\in M_m$ is
\[
    \left.\frac{\mathrm d}{\mathrm dt}\Psi(b+th)\right|_{t=0}
    =-\langle r-Kb,Kh\rangle.
\]
It vanishes for every $h\in M_m$ if and only if
$r-Kb\perp KM_m$.  For non-zero $h\in M_m$, injectivity and
non-negativity of $K$ imply $\langle Kh,h\rangle>0$ and
$\|Kh\|^2>0$.  These are the quadratic forms governing the response and
moment criteria, respectively, so both restrictions are strictly convex and
the two stationarity conditions are sufficient and uniquely determine their
minimisers.

For the response solution, take the inner product of
$r-K\sum_k a_kg_k$ with $g_j$.  Self-adjointness and
$g_j=K^{j-1}r$ give
\[
 \langle Kg_k,g_j\rangle
 =\langle K^kr,K^{j-1}r\rangle
 =\langle r,K^{j+k-1}r\rangle,
\]
which yields the first coefficient system.  For the moment solution, take
the inner product of the same residual with $Kg_j=K^jr$.  Then
\[
 \langle Kg_k,Kg_j\rangle
 =\langle K^kr,K^jr\rangle
 =\langle r,K^{j+k}r\rangle,
\]
which yields the second system.  The one-power shift proves that equality is
not implied merely by using the same $M_m$.
\end{proof}

\begin{proposition}[Finite termination with conjugate directions]
\label{prop:app-conjugate-direction-termination}
Let $A\in\mathbb R^{d\times d}$ be symmetric positive definite and let
$p_0,\ldots,p_{d-1}$ be non-zero, pairwise $A$-conjugate directions.  Starting
from any $x_0$, successively minimise
$f(x)=\frac12x^\top Ax-b^\top x$ along these directions without changing
the coefficients already selected.  After at most $d$ directions, the
iterate is the unique minimiser $A^{-1}b$.  Standard conjugate gradient
constructs such directions in exact arithmetic.
\end{proposition}

\begin{proof}
If $\sum_jc_jp_j=0$, multiplication by $p_k^\top A$ gives
$c_kp_k^\top Ap_k=0$.  Positive definiteness implies $c_k=0$ for every $k$,
so the $d$ directions are linearly independent and form a basis of
$\mathbb R^d$.  Write an iterate as
$x_0+\sum_j\eta_jp_j$.  Then
\[
 f\!\left(x_0+\sum_j\eta_jp_j\right)
 =f(x_0)+\sum_j\eta_jp_j^\top(Ax_0-b)
  +\frac12\sum_j\eta_j^2p_j^\top Ap_j,
\]
because all cross terms vanish by conjugacy.  The minimising coefficient for
one direction is therefore independent of every other coefficient.  After
all basis directions have been minimised, every directional derivative is
zero; hence $Ax-b=0$ and $x=A^{-1}b$.  The standard CG recurrence produces
pairwise $A$-conjugate search directions by induction from the residual
orthogonality relations, so the conclusion applies to its exact-arithmetic
iterates.  This result concerns the standard finite-dimensional CG geometry
used in Figure~\ref{fig:methods-cg-sd-geometry}; the released functional
recurrence is defined separately in \eqref{eq:methods-cg-alpha}--
\eqref{eq:methods-cg-direction-update}.
\end{proof}

\section{Late stopping and inference}
\label{sec:app-inference-proofs}

The asymptotic arguments in this section follow the inference construction of
\textcite{babii_carrasco_tsafack_2026}, with the identification qualification
made explicit.

The sampling operator retains the notation $T_n$ used throughout the
thesis.  To avoid confusing that operator with the scalar test statistic,
the latter is denoted by $\mathcal T_n$.  If $K$ is not injective, the
statistically identifiable null is $H_0:K\beta=Kb$; this becomes the literal
slope null $H_0:\beta=b$ when $\beta$ and $b$ lie in
$\mathcal H_{\mathrm{id}}$.  Chapter~\ref{chap:inference-methodology} writes
$\mathcal T_{n,m}$ when the finite iteration must be visible and abbreviates
$\mathcal T_{n,m_n}$ to $\mathcal T_n$ for an inference-specific index
sequence.

\begin{proposition}[Finite-iteration decomposition]
\label{prop:app-inference-decomposition}
For any candidate slope $b$ and any iterate $\widehat\beta_m$, define
\[
 u_n(b)=\widehat r-\widehat Kb,
 \qquad
 e_m=\widehat r-\widehat K\widehat\beta_m.
\]
Then
\[
 \widehat K(\widehat\beta_m-b)=u_n(b)-e_m.
\]
\end{proposition}

\begin{proof}
Add and subtract $\widehat r$:
\[
 \widehat K(\widehat\beta_m-b)
 =\widehat K\widehat\beta_m-\widehat Kb
 =(\widehat r-\widehat Kb)
  -(\widehat r-\widehat K\widehat\beta_m).
\]
\end{proof}

\begin{theorem}[Late-stopping equivalence and null limit]
\label{thm:app-late-stopping}
Suppose that $(X_i,Y_i)_{i=1}^n$ are independent and identically distributed
and satisfy
\[
    Y_i^c=\langle X_i^c,\beta\rangle+\varepsilon_i,
    \qquad \mathbb E(\varepsilon_i\mid X_i)=0.
\]
Assume $\mathbb E\|X^c\|^4<\infty$ and
$\mathbb E(\varepsilon^2\mid X)\leq\sigma_{\max}^2<\infty$ almost surely,
and let the identifiable null
$H_0:K\beta=Kb$ hold.  Equivalently, if $\beta,b\in
\mathcal H_{\mathrm{id}}$, assume $H_0:\beta=b$.  If, for an index
$m=m_n$ that may depend on $n$,
\[
 \sqrt n\,\|e_m\|=o_{\mathbb P}(1),
\]
then
\[
 \mathcal T_n(b)
 =n\|\widehat K(\widehat\beta_m-b)\|^2
 =n\|u_n(b)\|^2+o_{\mathbb P}(1).
\]
Let
\[
    V=\mathbb E\{\varepsilon^2(X^c\otimes X^c)\}
\]
be the covariance operator of $\varepsilon X^c$.  If its positive
eigenvalues are $\{\omega_j:j\in\mathcal J_V\}$, where $\mathcal J_V$ is
finite or countably infinite, then
\[
 \mathcal T_n(b)\rightsquigarrow
 \sum_{j\in\mathcal J_V}\omega_jZ_j^2,
 \qquad Z_j\stackrel{\mathrm{iid}}{\sim}\mathcal N(0,1).
\]
The sum is finite when $V$ has finite rank and converges almost surely and
in $L^1$ otherwise.
\end{theorem}

\begin{proof}
Under $H_0:K\beta=Kb$, Proposition~\ref{prop:app-identification} implies
$\langle X^c,\beta-b\rangle=0$ almost surely.  With population centring,
\[
 u_n(b)=\frac1n\sum_{i=1}^n\varepsilon_iX_i^c.
\]
If sample means are used instead, the exact expression is
\[
 u_n(b)=\frac1n\sum_{i=1}^n\varepsilon_iX_i^c
        -\overline\varepsilon\,\overline X^c.
\]
The two sample means are $O_{\mathbb P}(n^{-1/2})$ under the stated second
moments, so their product is $O_{\mathbb P}(n^{-1})$ and is negligible after
multiplication by $\sqrt n$.

The random element $\varepsilon X^c$ is mean zero because
$\mathbb E(\varepsilon\mid X)=0$, and it has a finite second moment by
the conditional variance bound and $\mathbb E\|X^c\|^2<\infty$.  The central
limit theorem for independent, identically
distributed random elements in a separable Hilbert space therefore gives
$\sqrt n\,u_n(b)\rightsquigarrow G$, where $G$ is a mean-zero Gaussian
element with covariance operator $V$.  In particular,
$\sqrt n\|u_n(b)\|=O_{\mathbb P}(1)$.

By Proposition~\ref{prop:app-inference-decomposition},
\begin{align*}
 \left|\mathcal T_n(b)-n\|u_n(b)\|^2\right|
 &\leq 2\sqrt n\|u_n(b)\|\sqrt n\|e_m\|
       +n\|e_m\|^2 \\
 &=o_{\mathbb P}(1).
\end{align*}
The continuous mapping theorem now gives
$n\|u_n(b)\|^2\rightsquigarrow\|G\|^2$.  Applying the covariance spectral
representation to $G$ yields
$G=\sum_j\sqrt{\omega_j}Z_jw_j$ in mean square for orthonormal eigenfunctions
$w_j$.  Parseval's identity therefore gives
$\|G\|^2=\sum_j\omega_jZ_j^2$.  Since
$\sum_j\omega_j=\operatorname{tr}(V)
=\mathbb E\|\varepsilon X^c\|^2<\infty$, the non-negative series converges
almost surely and in $L^1$, completing the proof.
\end{proof}

\begin{proposition}[Homoskedastic form of the variance operator]
\label{prop:app-homoskedastic-variance}
If $\mathbb E(\varepsilon^2\mid X)=\sigma^2$ almost surely, then the variance
operator in Theorem~\ref{thm:app-late-stopping} satisfies
\[
    V=\sigma^2K.
\]
Consequently, if $Kv_j=\lambda_jv_j$, then
$Vv_j=\sigma^2\lambda_jv_j$ and the null limit is
$\sum_j\sigma^2\lambda_jZ_j^2$ over the positive spectral indices.
\end{proposition}

\begin{proof}
For arbitrary $g,h\in\mathcal H$, iterated expectation gives
\begin{align*}
    \langle Vh,g\rangle
    &=\mathbb E\!\left[
       \varepsilon^2\langle X^c,h\rangle\langle X^c,g\rangle
      \right] \\
    &=\mathbb E\!\left[
       \mathbb E(\varepsilon^2\mid X)
       \langle X^c,h\rangle\langle X^c,g\rangle
      \right] \\
    &=\sigma^2\langle Kh,g\rangle.
\end{align*}
Equality of these bilinear forms for every $g$ and $h$ proves
$V=\sigma^2K$.  The spectral conclusion follows immediately.
\end{proof}

\begin{corollary}[Asymptotic size]
\label{cor:app-asymptotic-size}
Suppose the conditions of Theorem~\ref{thm:app-late-stopping} hold and $V$ has
at least one positive eigenvalue.  Let $q_{1-\alpha}$ be a continuity-point
$(1-\alpha)$ quantile of $Q=\sum_j\omega_jZ_j^2$.  Then the rejection rule
\[
    \varphi_n(b)=\mathbf 1\{\mathcal T_n(b)>q_{1-\alpha}\}
\]
has limiting null rejection probability $\alpha$.
\end{corollary}

\begin{proof}
Theorem~\ref{thm:app-late-stopping} gives
$\mathcal T_n(b)\rightsquigarrow Q$.  Since at least one weight is positive,
$Q$ has a continuous distribution: condition on all terms except one positive
weighted chi-squared term and integrate its continuous distribution function.
The portmanteau theorem at $q_{1-\alpha}$ therefore yields
\[
    \mathbb P\{\mathcal T_n(b)>q_{1-\alpha}\}
    \longrightarrow
    \mathbb P\{Q>q_{1-\alpha}\}=\alpha.
\]
\end{proof}

\begin{proposition}[Inference under fixed and local alternatives]
\label{prop:app-inference-alternatives}
Assume the conditions needed for the Hilbert-space central limit theorem in
Theorem~\ref{thm:app-late-stopping},
$\|\widehat K-K\|_{\mathrm{op}}=o_{\mathbb P}(1)$, and
$\sqrt n\|e_{m_n}\|=o_{\mathbb P}(1)$.

For a fixed alternative satisfying $K(\beta-b)\neq0$,
\[
    n^{-1}\mathcal T_n(b)
    \xrightarrow{\mathbb P}
    \|K(\beta-b)\|^2,
\]
so $\mathcal T_n(b)\to\infty$ in probability.  For the triangular sequence
$\beta_n=b+n^{-1/2}\Delta$ with fixed $\Delta\in\mathcal H$,
\[
    \mathcal T_n(b)\rightsquigarrow\|G+K\Delta\|^2.
\]
If $V\phi_j=\omega_j\phi_j$, the latter variable equals
\[
 \sum_j\omega_jZ_j^2
 +2\sum_j\sqrt{\omega_j}Z_j\langle\phi_j,K\Delta\rangle
 +\|K\Delta\|^2.
\]
\end{proposition}

\begin{proof}
Under a fixed alternative, the empirical null moment decomposes as
\[
 u_n(b)
 =\frac1n\sum_{i=1}^n\varepsilon_iX_i^c
  +\widehat K(\beta-b),
\]
apart from the lower-order sample-centring term treated in
Theorem~\ref{thm:app-late-stopping}.  The first term is
$O_{\mathbb P}(n^{-1/2})$, while operator convergence gives
$\widehat K(\beta-b)\to_{\mathbb P}K(\beta-b)$.  Since
$e_{m_n}=o_{\mathbb P}(n^{-1/2})$, Proposition
\ref{prop:app-inference-decomposition} implies
\[
 \widehat K(\widehat\beta_{m_n}-b)
 \xrightarrow{\mathbb P}K(\beta-b).
\]
Squaring the norm proves the first claim.

Under $\beta_n=b+n^{-1/2}\Delta$, multiply the analogous decomposition by
$\sqrt n$ to obtain
\[
 \sqrt n\,\widehat K(\widehat\beta_{m_n}-b)
 =\frac1{\sqrt n}\sum_{i=1}^n\varepsilon_iX_i^c
  +\widehat K\Delta-\sqrt n e_{m_n}+o_{\mathbb P}(1).
\]
The three displayed terms converge jointly to $G$, $K\Delta$, and zero,
respectively.  Slutsky's theorem and the continuous mapping theorem give the
local limit.  Finally write
$G=\sum_j\sqrt{\omega_j}Z_j\phi_j$ and expand
$\|G+K\Delta\|^2$; Parseval's identity justifies the displayed series.
\end{proof}

\begin{corollary}[Confidence sets by test inversion]
\label{cor:app-confidence-inversion}
Let $q_{1-\alpha}$ be a critical value for which the test that rejects when
$\mathcal T_n(b)>q_{1-\alpha}$ has asymptotic size $\alpha$ on an identified
parameter space $\Theta\subseteq\mathcal H_{\mathrm{id}}$.  Then
\[
 \mathcal C_{1-\alpha}
 =\{b\in\Theta:\mathcal T_n(b)\leq q_{1-\alpha}\}
\]
is an asymptotic $(1-\alpha)$ confidence set for every $\beta\in\Theta$.
\end{corollary}

\begin{proof}
By construction,
\[
 \mathbb P_\beta\{\beta\in\mathcal C_{1-\alpha}\}
 =\mathbb P_\beta\{\mathcal T_n(\beta)\leq q_{1-\alpha}\}
 =1-\mathbb P_\beta\{\mathcal T_n(\beta)>q_{1-\alpha}\}.
\]
The final rejection probability converges to $\alpha$, so the coverage
probability converges to $1-\alpha$.
\end{proof}

\begin{proposition}[Finite-dimensional confidence ellipsoid and envelope]
\label{prop:app-confidence-ellipsoid}
Let $H\in\mathbb R^{T\times p}$ contain the values of $p$ candidate basis
functions on an equally weighted grid, set $A=\widehat KH$, and let
$f=\widehat K\widehat\beta_m$.  For a critical value $q>0$, define
\[
    C_q=\left\{c\in\mathbb R^p:
      \frac nT\|f-Ac\|_2^2\leq q\right\}.
\]
If $A$ has full column rank, put
\[
    \widehat c=(A^\top A)^{-1}A^\top f,
    \qquad
    \rho^2=\frac{Tq}{n}-\|f-A\widehat c\|_2^2.
\]
Then $C_q$ is empty when $\rho^2<0$ and, when $\rho^2\geq0$,
\[
    C_q=\{c:(c-\widehat c)^\top A^\top A(c-\widehat c)
                    \leq\rho^2\}.
\]
For an evaluation vector $h(s)=(h_1(s),\ldots,h_p(s))^\top$, the extrema of
$h(s)^\top c$ over a non-empty $C_q$ are
\[
    h(s)^\top\widehat c
    \mathbin{\pm}
    \sqrt{\rho^2h(s)^\top(A^\top A)^{-1}h(s)}.
\]
If $A$ is rank deficient and $C_q$ is non-empty, its coefficient set is
unbounded along $\operatorname{Null}(A)$; a pointwise functional envelope is
unbounded wherever evaluation does not annihilate those directions.
\end{proposition}

\begin{proof}
The least-squares residual $u=f-A\widehat c$ is orthogonal to
$\operatorname{Range}(A)$.  Therefore, for every $c$,
\begin{align*}
    \|f-Ac\|_2^2
    &=\|u+A(\widehat c-c)\|_2^2 \\
    &=\|u\|_2^2
      +(c-\widehat c)^\top A^\top A(c-\widehat c).
\end{align*}
Substitution into the defining inequality proves the empty-set and ellipsoid
claims.  When $A$ has full column rank, write
$d=(A^\top A)^{1/2}(c-\widehat c)$.  The constraint becomes
$\|d\|_2\leq\rho$, and
\[
 h(s)^\top(c-\widehat c)
 =\left\{(A^\top A)^{-1/2}h(s)\right\}^\top d.
\]
Cauchy--Schwarz gives the displayed upper and lower values, attained when $d$
is parallel or antiparallel to $(A^\top A)^{-1/2}h(s)$.

If $A$ is rank deficient, take non-zero $v\in\operatorname{Null}(A)$.  For
any accepted $c$, every $c+tv$ has the same residual for all $t\in\mathbb R$.
Thus the coefficient set is unbounded, and so is the value
$h(s)^\top(c+tv)$ whenever $h(s)^\top v\neq0$.
\end{proof}

Late stopping removes iteration bias. Identifying a unique slope still
requires restriction to $\operatorname{Null}(K)^\perp$.
